\documentclass[a4paper,10pt]{report}\input{../0.Préambule/Préambule_cours_prof}\begin{document}

\newpage
\section*{Calcul littéral 2}
\addcontentsline{toc}{section}{Calcul littéral 2}

\subsection*{Développement}
\addcontentsline{toc}{subsection}{Développement}

\begin{exer1}[Produit ?]
\begin{enumerate}[font=\color{exer1}]
\item Souligne ci-dessous les expressions qui sont des produits et entoure leurs facteurs
\begin{enumerate}[font=\color{exer1}]
\begin{tabularx}{\linewidth}{*{4}{>{\arraybackslash}X}}
\item $P = 5 \times x - 4$ &
\item $R = 5 \times (a - 4)$ &
\item $A = 4y \times (-3y)$ &
\item $L = 5[2x+6]$ \\
\item $I = (-2+x) \times 5x$ &
\item $N = 3u + 2(u-5)$ &
\item $E = (3x+2)(x-5)$ &
\item $S = 3v + 2 \times v - 5$ \\
\end{tabularx}
\end{enumerate}
\item Parmi les expressions précédentes, lesquelles pourrais-tu développer ?
\end{enumerate}
\end{exer1}

\begin{exer1}
Développe les expressions ci-dessous (On ne vous demande pas de faire les calculs !!)
\begin{enumerate}[font=\color{exer1}]
\begin{tabularx}{\linewidth}{*{2}{>{\arraybackslash}X}}
\item $L = 36 \times (21 + 55) = \hspace{.8cm}  \times \hspace{.8cm} + \hspace{.8cm} \times $ &
\item $O = 81 \times (48 - 7) = \hspace{.8cm}  \times \hspace{.8cm} - \hspace{.8cm} \times $ \\
\item $U = (85-7) \times 71 = $ &
\item $P = (32 + 91) \times 44 = $ \\
\end{tabularx}
\end{enumerate}
\end{exer1}

\begin{exer1}

\begin{enumerate}[font=\color{exer1}]
\item Complète le tableau suivant.

\vspace{-8mm}
\begin{flushright}
\renewcommand{\arraystretch}{1.5}
\begin{tabularx}{6.5cm}{|>{\columncolor{exer1!20}\centering}p{1.3cm}|*{3}{>{\centering \arraybackslash}X|}}
\hline
\rowcolor{exer1!20} $\times$ & $100$ & $1$ & $2$ \\\hline
$24$ & & &  \\\hline
\end{tabularx}
\end{flushright}

\vspace{-8mm}
\item Donne, sans calculatrice, le résultat des produits suivants
\begin{enumerate}[font=\color{exer1}]
\begin{tabularx}{\linewidth}{*{4}{>{\arraybackslash}X}}
\item $24 \times 101$ &
\item $24 \times 99$ &
\item $24 \times 102$ &
\item $24 \times 98$ \\
\end{tabularx}
\end{enumerate}
\end{enumerate}
\end{exer1}

\begin{exer1}
Développe les expressions suivantes :
\begin{enumerate}[font=\color{exer1}]
\begin{tabularx}{\linewidth}{*{5}{>{\arraybackslash}X}}
\item $A=3(x+6)$ &
\item $B=5(6-y)$ &
\item $I=-8(-5-3y)$ &
\item $O=6(4x-9)$ &
\item $G=-12(-5+3z)$ \\
\item $L=(-3+y) \times 9$ &
\item $Y=-6(2x-7)$ &
\item $P=(3t+2) \times 8$ &
\item $H=-8(9-7x)$ &
\item $E=-8z(4-3z)$ \\
\end{tabularx}
\end{enumerate}
\end{exer1}

\begin{exer2}
On considère l'expression $A= 3x + 5x(x-2)$
\begin{enumerate}[font=\color{exer2}]
\item Souligne l'opération prioritaire
\item Développe et réduis l'expression $A$
\end{enumerate}
\end{exer2}

\begin{exer2}
On considère l'expression $B = 4 - 2(3-5u)$
\begin{enumerate}[font=\color{exer2}]
\item Complète : $B = 4 + ( \hspace{8mm} ) \times (3-5u)$
\item Développe et réduis l'expression $B$
\end{enumerate}
\end{exer2}

\begin{exer2}
On considère l'expression $C = 3x - (2x+5) \times 4$
\begin{enumerate}[font=\color{exer2}]
\item Ajoute des crochets autour de l'opération prioritaire.
\item Développe et réduis l'expression $C$
\end{enumerate}
\end{exer2}

\begin{exer2}
Développe et réduis les expressions suivantes :
\begin{enumerate}[font=\color{exer2}]
\begin{tabularx}{\linewidth}{*{4}{>{\arraybackslash}X}}
\item $Y=3(x+6)-2$ &
\item $U=4+3(2y-2)$ &
\item $K=3,5(2-x)+8,2$ &
\item $A=9(t-6)+2t$ \\
\end{tabularx}
\end{enumerate}
\end{exer2}

\begin{exer2}
Développe et réduis les expressions suivantes :
\begin{enumerate}[font=\color{exer2}]
\begin{tabularx}{\linewidth}{*{4}{>{\arraybackslash}X}}
\item $N=x(x+6)-x$ &
\item $O=x(y-2)+xy$ &
\item $I=2x(x+1)-4x^2$ &
\item $X=3x(x+4)-6x^2$ \\
\end{tabularx}
\end{enumerate}
\end{exer2}

\begin{exer2}
Développe et réduis les expressions suivantes :
\begin{enumerate}[font=\color{exer2}]
\begin{tabularx}{\linewidth}{*{3}{>{\arraybackslash}X}}
\item $P=-2(x-5)-3(7-4x)$ &
\item $O=8+2y-5(2y-6)+4$ &
\item $N=3x+5x(4-2x)-2(x^2-3x+5)$ \\
\item $E=8-2x(3x-4)+5x(3-x)$ &
\item $Y=7y+4(3y-6)+3+2(3y-7)$ &\\
\end{tabularx}
\end{enumerate}
\end{exer2}

\newpage
\subsection*{Applications}
\addcontentsline{toc}{subsection}{Applications}

\begin{exer3}
On considère le programme de calcul ci-contre :

\begin{enumerate}[font=\color{exer3}]
\item Exécute ce programme pour $x=3$, puis pour $x=2$
\item Que remarques-tu ?
\item Quelle expression obtiens tu si le nombre départ est $x$ ?
\item Explique alors ta réponse à la question \textbf{2.}
\end{enumerate}

\vspace{-25mm}
\begin{flushright}
\arrayrulecolor{black}
\renewcommand{\arraystretch}{1.2}
\begin{tabular}{|l|}
\hline
\edupythoninline{Choisis un nombre}\\
\edupythoninline{Soustrais 8 a ce nombre}\\
\edupythoninline{Multiplie le resultat par -4}\\
\edupythoninline{Ajoute le quadruple du nombre de depart}\\
\hline
\end{tabular}
\end{flushright}
\end{exer3}

\begin{exer3}

\begin{enumerate}[font=\color{exer3}]
\item Applique le programme de calcul suivant pour 

\noindent deux valeurs de ton choix :

\vspace{-12mm}
\begin{flushright}
\arrayrulecolor{black}
\renewcommand{\arraystretch}{1.2}
\begin{tabular}{|l|}
\hline
\edupythoninline{Choisis un nombre}\\
\edupythoninline{Soustrais le a 5}\\
\edupythoninline{Multiplie le resultat par 4}\\
\edupythoninline{Ajoute le triple du nombre de depart}\\
\hline
\end{tabular}
\end{flushright}

\vspace{-12mm}
\item Marlène dit que ce programme pourrait ne contenir \\que deux instructions au lieu de quatre. Lesquelles ?
\end{enumerate}
\end{exer3}

\begin{exer3}
Montre que les deux solides ci-dessous ont le même volume.

\begin{center}
\includegraphics[scale=0.35]{../40.Image/11. Exo-Calcul_littéral_2_1.png}
\end{center}
\end{exer3}

\begin{exer3}
Dans un parc zoologique, les enfants paient $3$ \euro de moins que les adultes. On appelle $p$ le prix d'entrée d'un enfant. Aujourd'hui, $130$ adultes et $140$ enfants sont venus au zoo.

\begin{enumerate}[font=\color{exer3}]
\item Exprime en fonction de $p$ la recette réalisée par le zoo aujourd'hui.
\item Quelle est la recette si le ticket adulte coute $15$\euro ?
\end{enumerate}
\end{exer3}

\begin{exer3}
On considère le programme de calcul suivant :

\vspace{-12mm}
\begin{flushright}
\arrayrulecolor{black}
\renewcommand{\arraystretch}{1.2}
\begin{tabular}{|l|}
\hline
\edupythoninline{Choisis un nombre}\\
\edupythoninline{Augmenter le nombre de 5}\\
\edupythoninline{Multiplie le resultat par 4}\\
\edupythoninline{Oter le quadruple du nombre de depart}\\
\edupythoninline{Oter 10 et annoncer le resultat}\\
\hline
\end{tabular}
\end{flushright}

\vspace{-20mm}
\begin{enumerate}[font=\color{exer3}]
\item Applique ce programme de calcul à $5$ et à $2,3$.
\item Que remarques-tu ?
\item Pour chaque étape du programme, complète le diagramme \\ par des expressions simplifiées.

\includegraphics[scale=0.35]{../40.Image/11. Exo-Calcul_littéral_2_2.png}
\item Conclus
\end{enumerate}
\end{exer3}

\begin{exer3}
Voici un programme :

\vspace{-1cm}
\noindent \begin{minipage}{9.5cm}
\begin{enumerate}[font=\color{exer3}]
\item Détermine le résultat obtenu pour le nombre de ton choix.
\item Exprime le résultat obtenu par le programme pour un nombre $x$
\item Lisa remarque qu'en choisissant un nombre entier, le programme donne toujours un multiple de $3$. Justifie cette remarque.
\end{enumerate}
\end{minipage}
\hspace{5mm}
\begin{minipage}{7cm}
\begin{scratch}[]
\blockinit{quand \greenflag est cliqué}
\blockmove{demander \ovalnum{Choisis un nombre} et attendre}
\blockvariable{mettre {\selectmenu{A} à \ovaloperator{\ovaloperator{\ovalnum{3} * \ovalsensing{réponse}}+\ovalnum{7}}}}
\blockvariable{mettre {\selectmenu{Résultat} à \ovaloperator{\ovalvariable{Résultat} * \ovalsensing{réponse}}}}
\blockvariable{mettre {\selectmenu{Résultat} à \ovaloperator{\ovalvariable{Résultat} - \ovalsensing{réponse}}}}
\blocklook{dire \ovaloperator{regrouper \ovalnum{Le résultat est } et \ovalvariable{Résultat}}}
\end{scratch}
\end{minipage}
\end{exer3}

\end{document}
